\section{What the Data Says About India's Overnight Liquidity}
\label{sec:results}

This section interprets the headline findings from
Section~\ref{sec:methods} as substantive statements about the current
state of India's overnight call-money market. Five claims follow,
each grounded directly in a specific number or figure from the
analysis.

\subsection{The 2020 break is structural, not transitional}

The K-Means clustering with $K=2$ partitions the 545-week panel
cleanly into a pre-2020 segment of 315 weeks and a post-2020 segment
of 230 weeks (Section~\ref{sec:methods-regimes},
Figure~\ref{fig:regime-state-machine}). The cluster boundary is not
gradual: the silhouette score drops sharply once $K$ goes above 2,
and the PCA projection in Figure~\ref{fig:pca-scatter} shows two
clusters separated along PC1 with a narrow transition band only at
the COVID weeks themselves. This is empirical evidence that the post-
2020 monetary regime is structurally distinct from the pre-2020
regime, and that even the 2022--23 hiking cycle continues to operate
inside the new framework rather than reverting to the old one.

The 1.72 percentage-point gap in mean WACMR between the two regimes
(6.528\% versus 4.809\%) is not a transient feature of any one year.
Within the post-2020 regime, mean WACMR rose from 3.5\% in 2020 to
6.7\% in early 2024, but the cluster assignment did not flip.
Mechanically, what changed in 2020 was the introduction of Variable
Rate Repo and Variable Rate Reverse Repo as RBI's primary liquidity
tools; the persistence of the regime label through 2024 reflects the
fact that these tools remain in active use, and that the operational
playbook has not reverted.

\subsection{The corridor explains almost everything; equity and forex
explain almost nothing}

The SHAP analysis (Table~\ref{tab:shap-top15}) is unambiguous on
this point. The top five drivers of weekly WACMR forecasts are
\texttt{target\_lag1} (mean $|$SHAP$|$ 0.48959), \texttt{repo\_lag1}
(0.21053), the current Repo Rate (0.19459), the engineered
\texttt{spread\_wacmr\_minus\_repo} (0.07247), and the MSF Rate
(0.05847). These five features are either direct policy-corridor rates
or short-lag autoregressive functions of WACMR itself. Their combined
mean absolute SHAP value of 1.02565 dwarfs the contribution of every
other feature.

By contrast, the twenty-eight equity and currency features added
through the supplemental Yahoo Finance pipeline (Nifty 50 OHLCV,
USD/INR OHLCV, plus five custom technical indicators per instrument)
do not appear in the top 15. Their cumulative contribution is below
the contribution of any single corridor feature. The first market
feature to appear, \texttt{usdinr\_ImpulseMACD}, sits at rank 23 with
mean $|$SHAP$|$ of 0.0038, which is two orders of magnitude smaller
than the headline corridor features.

\begin{finding}
At the weekly horizon, India's call-money market is overwhelmingly
LAF-corridor-bound. WACMR moves with the announced policy rate and
its own recent history; it does not respond meaningfully to equity
momentum, currency momentum, or any combination of technical indicators
on those instruments. This is the central empirical claim of the
report.
\end{finding}

\subsection{Forecasting is feasible to about ten basis points at a
one-week horizon, but explicit regime tagging adds nothing}

The walk-forward XGBoost achieves RMSE 0.1019 (about ten basis points
on the WACMR scale) and 70.9\% directional accuracy across 389
out-of-sample test weeks (Table~\ref{tab:perf}). Both numbers are
high enough to be operationally useful for treasury hedging and
overnight-fund positioning. RMSE of 0.1019 corresponds to a 95\%
prediction interval of roughly $\pm 20$ basis points around the point
forecast.

The regime-aware variant, with two one-hot regime indicators and two
PCA-space distance features added, performs marginally worse on RMSE
(0.1044, $+2.5$\%) and identically on MAE and directional accuracy.
This is a clean empirical refutation of hypothesis H2. The natural
reading is that the lagged Repo Rate and the autoregressive lags of
WACMR already encode whatever the K-Means regime label encodes;
adding the explicit label introduces correlated noise without new
information. The implication for downstream practitioners is that
investing engineering effort in regime-tagging machinery is not the
highest-value use of time. A simple feature set, kept clean, performs
as well as a more elaborate one.

\subsection{Tail errors come from policy discontinuities, not
calendar pressure}

The residual analysis in Figure~\ref{fig:residual-calendar} contains
one of the more counterintuitive findings of the project. The
calendar-event hypothesis, that financial-year end and the four
advance-tax dates produce the model's worst weeks, is rejected. RMSE
on the 66 calendar-event weeks is 0.0870, which is meaningfully
better than the 0.1046 RMSE on the 323 non-event weeks. Ten years of
training data is enough for the model to learn the seasonal pattern.

The actual tail residuals concentrate on three families of events
that have no historical analogue in training data:

\begin{itemize}
\item COVID emergency rate cuts in March--April 2020 and the COVID
  second-wave peak in May 2021. The April 2020 prediction missed by
  $-0.66$ percentage points, the largest single error in the
  out-of-sample set.
\item The surprise 40 basis-point inter-meeting Repo hike on 13 May
  2022, which the model missed by $-0.42$ percentage points.
\item The post-hike re-tightening shocks in 2022 (June, September,
  October), which the model missed by 0.30 to 0.34 percentage points.
\end{itemize}

These are by definition not forecastable from historical features
alone. They are unsignalled discretionary policy actions. Any
forecasting effort that wants to handle them needs to inject an
exogenous policy-action prior, for example a decision-day dummy
based on the published MPC calendar.

\subsection{The WACMR-Repo spread is a real-time liquidity barometer}

The fourth-ranked SHAP feature is the engineered
\texttt{spread\_wacmr\_minus\_repo}, with mean absolute SHAP 0.0725.
The mean spread by regime tells a clear story: in Regime 1
(Tightening) the spread averages roughly zero, whereas in Regime 0
(Accommodation) it averages $-0.19$ percentage points, with WACMR
sitting consistently below the policy Repo Rate. When the spread
flips meaningfully positive (WACMR > Repo) for several consecutive
weeks, it has historically preceded liquidity-deficit interventions.
This is not a new observation in the academic literature, but the
SHAP analysis quantifies the predictive weight: the spread is more
informative than the MSF Rate itself.

\subsection{Summary of evidence}

Table~\ref{tab:hypothesis-results} summarises the verdict on each
hypothesis stated in Section~\ref{sec:problem}.

\begin{table}[H]
\caption{Hypothesis verdicts at the end of the analysis.}
\label{tab:hypothesis-results}
\renewcommand{\arraystretch}{1.2}
\rowcolors{2}{SoftGray}{white}
\begin{tabularx}{\linewidth}{@{} l c Y @{}}
\toprule
\rowcolor{Slate}\color{white}\textbf{Hypothesis} &
\color{white}\textbf{Verdict} &
\color{white}\textbf{Evidence} \\
\midrule
H1, regime structure exists & Supported & K-Means $K=2$, silhouette 0.4238, two non-overlapping regimes aligned with COVID. \\
H2, regime labels improve forecasts & Rejected & Regime-aware RMSE 0.1044 vs baseline 0.1019 ($+2.5$\%); MAE and directional accuracy unchanged. \\
H3, corridor rates dominate & Supported & Repo, Reverse Repo, and MSF together with the autoregressive lags occupy the entire SHAP top 15. \\
H4, equity and forex add value & Rejected & None of 28 supplemental features cracks the SHAP top 15; first market feature ranks 23 with cumulative impact two orders of magnitude below the corridor. \\
\bottomrule
\end{tabularx}
\end{table}
